
\documentclass[12pt,a4paper]{article}

%--------------------------------------------------------------
% PACKAGES
%--------------------------------------------------------------
\usepackage[margin=1in, headheight=14pt]{geometry}
\usepackage{amsmath, amssymb}
\usepackage{booktabs}
\usepackage{array}
\usepackage{longtable}
\usepackage{xcolor}
\usepackage{listings}
\usepackage{graphicx}
\usepackage{hyperref}
\usepackage{parskip}
\usepackage{titlesec}
\usepackage{enumitem}
\usepackage{fancyhdr}
\usepackage{float}
\usepackage{caption}
\usepackage{multirow}
\usepackage{tabularx}
\usepackage[T1]{fontenc}

%--------------------------------------------------------------
% COLOURS
%--------------------------------------------------------------
\definecolor{codebg}{RGB}{245,245,245}
\definecolor{codeframe}{RGB}{180,180,180}
\definecolor{kwcolor}{RGB}{0,0,180}
\definecolor{strcolor}{RGB}{163,21,21}
\definecolor{cmtcolor}{RGB}{0,128,0}
\definecolor{titleblue}{RGB}{0,70,127}
\definecolor{sectionblue}{RGB}{31,73,125}
\definecolor{outputbg}{RGB}{240,248,255}
\definecolor{outputframe}{RGB}{100,149,237}
\definecolor{conclusionbg}{RGB}{255,250,240}
\definecolor{conclusionframe}{RGB}{210,150,50}

%--------------------------------------------------------------
% LISTINGS STYLE (Python)
%--------------------------------------------------------------
\lstdefinestyle{pythonstyle}{
    backgroundcolor=\color{codebg},
    frame=single,
    rulecolor=\color{codeframe},
    basicstyle=\ttfamily\small,
    keywordstyle=\color{kwcolor}\bfseries,
    stringstyle=\color{strcolor},
    commentstyle=\color{cmtcolor}\itshape,
    numberstyle=\tiny\color{gray},
    numbers=left,
    numbersep=6pt,
    breaklines=true,
    breakatwhitespace=false,
    showstringspaces=false,
    tabsize=4,
    language=Python,
    captionpos=b,
    xleftmargin=15pt,
    xrightmargin=5pt,
    aboveskip=8pt,
    belowskip=4pt,
}

\lstdefinestyle{outputstyle}{
    backgroundcolor=\color{outputbg},
    frame=single,
    rulecolor=\color{outputframe},
    basicstyle=\ttfamily\small,
    breaklines=true,
    numbers=none,
    xleftmargin=10pt,
    xrightmargin=5pt,
    aboveskip=6pt,
    belowskip=4pt,
}

%--------------------------------------------------------------
% SECTION FORMATTING
%--------------------------------------------------------------
\titleformat{\section}[block]
  {\large\bfseries\color{sectionblue}}
  {\thesection.}{0.5em}{}[\titlerule]

\titleformat{\subsection}[block]
  {\normalsize\bfseries\color{sectionblue}}
  {\thesubsection}{0.5em}{}

\titleformat{\subsubsection}[block]
  {\normalsize\bfseries}
  {\thesubsubsection}{0.5em}{}

%--------------------------------------------------------------
% HEADER / FOOTER
%--------------------------------------------------------------
\pagestyle{fancy}
\fancyhf{}
\rhead{\small Assignment 1 -- Lung Cancer Data Analysis}
\lhead{\small Python \& Statistical Methods}
\rfoot{\small Page \thepage}
\lfoot{\small Biostatistics / Data Science}

%--------------------------------------------------------------
% HYPERREF SETUP
%--------------------------------------------------------------
\hypersetup{
    colorlinks=true,
    linkcolor=sectionblue,
    urlcolor=sectionblue,
    citecolor=sectionblue,
    pdftitle={Assignment 1 - Lung Cancer Data Analysis},
    pdfauthor={},
}

%--------------------------------------------------------------
% CUSTOM ENVIRONMENTS
%--------------------------------------------------------------
\newenvironment{mybox}[1]{%
  \par\medskip\noindent
  \textbf{#1}\par\smallskip
  \noindent\hrulefill\par\smallskip
  \noindent\ignorespaces
}{%
  \par\noindent\hrulefill\par\medskip
}

%--------------------------------------------------------------
% DOCUMENT
%--------------------------------------------------------------
\begin{document}

%==============================================================
% TITLE PAGE
%==============================================================
\begin{titlepage}
  \centering
  \vspace*{2cm}
  {\huge\bfseries\color{titleblue} Assignment 1 \par}
  \vspace{0.5cm}
  {\Large\bfseries Lung Cancer Data Analysis using Python \par}
  \vspace{1cm}
  \rule{0.8\linewidth}{1pt}\par
  \vspace{0.5cm}
  {\large Statistical Methods: Descriptive Statistics, Hypothesis Testing,\\
  Logistic Regression, ANOVA, and Count Models \par}
  \vspace{1cm}
  \rule{0.8\linewidth}{0.5pt}\par
  \vspace{1.5cm}
  {\normalsize
    \begin{tabular}{ll}
      \textbf{Dataset:} & Lung Cancer Survey Data ($n = 500$) \\[4pt]
      \textbf{Tools:}   & Python (pandas, scipy, statsmodels, sklearn) \\[4pt]
      \textbf{Date:}    & May 2026 \\
    \end{tabular}
  }
  \vfill
  {\small Prepared using \LaTeX{} from Python code, outputs, and interpretations.}
\end{titlepage}

\tableofcontents
\newpage

%==============================================================
\section{Problem 1: Lung Cancer Data Analysis}
%==============================================================

\textit{Considering the lung cancer dataset for solving Problem 1 using Python.}

%--------------------------------------------------------------
\subsection{(i) Distribution of Age -- Normality Analysis}
%--------------------------------------------------------------

\subsubsection{Python Code}

\begin{lstlisting}[style=pythonstyle, caption={Loading data and analysing Age distribution}]
import pandas as pd
import matplotlib.pyplot as plt
from scipy.stats import shapiro

df = pd.read_csv('himel.xlsx.csv')
print(df.head())

age = df['Age']

print("\nSummary Statistics of Age:")
print(age.describe())

print("\nSkewness:", age.skew())
print("Kurtosis:", age.kurt())

plt.hist(age, bins=10, edgecolor='black')
plt.title("Histogram of Age")
plt.xlabel("Age")
plt.ylabel("Frequency")
plt.show()

stat, p = shapiro(age)
print("\nShapiro-Wilk Test")
print("Statistic =", stat)
print("P-value =", p)

if p > 0.05:
    print("Age is approximately normally distributed.")
else:
    print("Age is not normally distributed.")
\end{lstlisting}

\subsubsection{Output}

\begin{lstlisting}[style=outputstyle, caption={Summary statistics and normality test output}]
Summary Statistics of Age:
count    500.000000
mean      43.904000
std       14.604841
min       20.000000
25%       30.750000
50%       45.000000
75%       56.000000
max       69.000000

Skewness:  0.012141455267137275
Kurtosis: -1.232106868020186

Shapiro-Wilk Test
Test Statistic: 0.9488992946952997
P-value:        4.0350550859741284e-12

Conclusion: Age is NOT normally distributed.
\end{lstlisting}

\subsubsection{Interpretation}

\begin{table}[H]
  \centering
  \caption{Descriptive Statistics for Age}
  \begin{tabular}{lc}
    \toprule
    \textbf{Statistic} & \textbf{Value} \\
    \midrule
    Count    & 500 \\
    Mean     & 43.90 \\
    Std Dev  & 14.60 \\
    Min      & 20 \\
    Q1       & 30.75 \\
    Median   & 45.00 \\
    Q3       & 56.00 \\
    Max      & 69 \\
    Skewness & 0.0121 \\
    Kurtosis & $-1.232$ \\
    \bottomrule
  \end{tabular}
\end{table}

\noindent\textbf{Skewness ($\approx 0.012$):} The \texttt{Age} variable is approximately symmetric.

\noindent\textbf{Kurtosis ($\approx -1.232$):} Platykurtic -- flatter than a normal distribution.

\noindent\textbf{Shapiro-Wilk Test:}
\begin{align*}
  W &= 0.9489, \quad p = 4.04 \times 10^{-12}
\end{align*}
Since $p \ll 0.05$, we \textbf{reject the null hypothesis} of normality. The \texttt{Age} variable is \textbf{approximately symmetric} but formally \textbf{not normally distributed}.

%--------------------------------------------------------------
\subsection{(ii) Proportion of Smokers vs.\ Non-Smokers}
%--------------------------------------------------------------

\subsubsection{Python Code}

\begin{lstlisting}[style=pythonstyle, caption={Smoking proportions}]
smoking_prop = df['Smoking'].value_counts(normalize=True) * 100
print(smoking_prop)
\end{lstlisting}

\subsubsection{Output}

\begin{lstlisting}[style=outputstyle]
No     58.6%
Yes    41.4%
\end{lstlisting}

\subsubsection{Interpretation}

\begin{table}[H]
  \centering
  \caption{Smoking Status Distribution}
  \begin{tabular}{lcc}
    \toprule
    \textbf{Category} & \textbf{Count} & \textbf{Percentage (\%)} \\
    \midrule
    Non-Smoker & 293 & 58.6 \\
    Smoker     & 207 & 41.4 \\
    \midrule
    Total      & 500 & 100.0 \\
    \bottomrule
  \end{tabular}
\end{table}

About \textbf{41.4\%} are smokers and \textbf{58.6\%} are non-smokers.

%--------------------------------------------------------------
\subsection{(iii) Income Group with Highest Frequency}
%--------------------------------------------------------------

\subsubsection{Python Code}

\begin{lstlisting}[style=pythonstyle, caption={Income group frequency}]
income_counts = df['Income'].value_counts()
print("Frequency of Each Income Group:")
print(income_counts)

highest_group = income_counts.idxmax()
highest_frequency = income_counts.max()
print(f"\nThe {highest_group} income group has the highest frequency "
      f"with {highest_frequency} individuals.")
\end{lstlisting}

\subsubsection{Output}

\begin{lstlisting}[style=outputstyle]
Frequency of Each Income Group:
Income
Low       204
Medium    197
High       99

The Low income group has the highest frequency with 204 individuals.
\end{lstlisting}

\subsubsection{Interpretation}

\begin{table}[H]
  \centering
  \caption{Income Group Frequencies}
  \begin{tabular}{lcc}
    \toprule
    \textbf{Income Group} & \textbf{Frequency} & \textbf{Percentage (\%)} \\
    \midrule
    Low    & 204 & 40.8 \\
    Medium & 197 & 39.4 \\
    High   &  99 & 19.8 \\
    \midrule
    Total  & 500 & 100.0 \\
    \bottomrule
  \end{tabular}
\end{table}

The \textbf{Low income group} has the highest frequency (40.8\%).

%--------------------------------------------------------------
\subsection{(iv) Percentage Exposed to High Pollution}
%--------------------------------------------------------------

\subsubsection{Python Code}

\begin{lstlisting}[style=pythonstyle, caption={Pollution exposure analysis}]
pollution_counts = df['Pollution'].value_counts()
print("Frequency of Pollution Exposure Categories:")
print(pollution_counts)

high_pollution_percentage = (df['Pollution'] == 'High').mean() * 100
print(f"\n{high_pollution_percentage:.2f}% of individuals are exposed to High pollution.")
\end{lstlisting}

\subsubsection{Output}

\begin{lstlisting}[style=outputstyle]
Pollution
High    352
Low     148

70.40% of individuals are exposed to High pollution.
\end{lstlisting}

\subsubsection{Interpretation}

\textbf{70.40\%} of individuals are exposed to high pollution, indicating the sample is drawn predominantly from high-pollution environments.

%--------------------------------------------------------------
\subsection{(v) Overall Prevalence of Lung Disease}
%--------------------------------------------------------------

\subsubsection{Python Code}

\begin{lstlisting}[style=pythonstyle, caption={Lung disease prevalence}]
lung_counts = df['LungDisease'].value_counts()
print("Frequency of Lung Disease Status:")
print(lung_counts)

prevalence = (df['LungDisease'] == 'Yes').mean() * 100
print(f"\n{prevalence:.2f}% of individuals have lung disease.")
\end{lstlisting}

\subsubsection{Output}

\begin{lstlisting}[style=outputstyle]
LungDisease
Yes    264
No     236

52.80% of individuals have lung disease.
\end{lstlisting}

\subsubsection{Interpretation}

The overall prevalence of lung disease is \textbf{52.80\%}, reflecting a high-risk study population.

%==============================================================
\section{Association and Crosstab Analysis}
%==============================================================

\subsection{Association Between Smoking and Lung Disease}

\subsubsection{Python Code}

\begin{lstlisting}[style=pythonstyle, caption={Chi-square test}]
from scipy.stats import chi2_contingency

table = pd.crosstab(df['Smoking'], df['LungDisease'])
print("Contingency Table:")
print(table)

chi2, p_value, dof, expected = chi2_contingency(table)
print("\nChi-Square Statistic =", chi2)
print("Degrees of Freedom =", dof)
print("P-value =", p_value)

if p_value < 0.05:
    print("Significant association between Smoking and Lung Disease.")
else:
    print("No significant association.")
\end{lstlisting}

\subsubsection{Output}

\begin{lstlisting}[style=outputstyle]
Contingency Table:
LungDisease    No  Yes
Smoking
No            184  109
Yes            52  155

Chi-Square Statistic = 67.5942649754021
Degrees of Freedom   = 1
P-value              = 2.0085094601569646e-16

Significant association between Smoking and Lung Disease.
\end{lstlisting}

\subsubsection{Row Percentage Table}

\begin{lstlisting}[style=pythonstyle]
contingency_table = pd.crosstab(
    df['Smoking'], df['LungDisease'],
    margins=True, normalize='index'
) * 100
print(contingency_table.round(2))
\end{lstlisting}

\begin{lstlisting}[style=outputstyle]
LungDisease     No    Yes
Smoking
No           62.80  37.20
Yes          25.12  74.88
All          47.20  52.80
\end{lstlisting}

\subsubsection{Interpretation}

\begin{table}[H]
  \centering
  \caption{Crosstab: Smoking vs.\ Lung Disease (Row \%)}
  \begin{tabular}{lccc}
    \toprule
    \textbf{Smoking} & \textbf{No Disease (\%)} & \textbf{Disease (\%)} & \textbf{Total} \\
    \midrule
    No  & 62.80 & 37.20 & 293 \\
    Yes & 25.12 & 74.88 & 207 \\
    \midrule
    All & 47.20 & 52.80 & 500 \\
    \bottomrule
  \end{tabular}
\end{table}

$\chi^2(1) = 67.59$, $p = 2.01 \times 10^{-16}$. Strong evidence of association. Among smokers, 74.88\% have lung disease vs.\ only 37.20\% of non-smokers.
PART 2 — Lines ~600 to end:

latex
%--------------------------------------------------------------
\subsection{Lung Disease Prevalence Across Income Groups}
%--------------------------------------------------------------

\subsubsection{Python Code}

\begin{lstlisting}[style=pythonstyle]
income_prevalence = pd.crosstab(
    df['Income'], df['LungDisease'], normalize='index'
) * 100
print(income_prevalence.round(2))
\end{lstlisting}

\subsubsection{Output}

\begin{lstlisting}[style=outputstyle]
LungDisease     No    Yes
Income
High         51.52  48.48
Low          44.12  55.88
Medium       48.22  51.78
\end{lstlisting}

\subsubsection{Interpretation}

\begin{table}[H]
  \centering
  \caption{Lung Disease Prevalence by Income Group}
  \begin{tabular}{lcc}
    \toprule
    \textbf{Income Group} & \textbf{No Disease (\%)} & \textbf{Disease (\%)} \\
    \midrule
    High   & 51.52 & 48.48 \\
    Medium & 48.22 & 51.78 \\
    Low    & 44.12 & 55.88 \\
    \bottomrule
  \end{tabular}
\end{table}

The \textbf{Low income group} has the highest lung disease prevalence (55.88\%), suggesting lower income is associated with a higher disease burden.

%--------------------------------------------------------------
\subsection{Lung Disease by Pollution Exposure}
%--------------------------------------------------------------

\begin{lstlisting}[style=outputstyle]
Lung Disease Prevalence by Pollution (%):
LungDisease     No    Yes
Pollution
High         38.64  61.36
Low          67.57  32.43

Correlation between Age and Lung Disease:     0.2497
Correlation between Smoking and Lung Disease: 0.3717
\end{lstlisting}

High-pollution individuals have a 61.36\% lung disease prevalence vs.\ 32.43\% for low-pollution. Smoking ($r=0.37$) has a stronger association than Age ($r=0.25$).

%==============================================================
\section{Statistical Testing}
%==============================================================

\subsection{(i) Chi-Square Test}

\begin{lstlisting}[style=outputstyle]
Chi-Square Statistic = 67.5942649754021
Degrees of Freedom   = 1
P-value              = 2.0085094601569646e-16

Conclusion: Statistically significant association between
Smoking and Lung Disease.
\end{lstlisting}

\subsection{(ii) Fisher's Exact Test}

\subsubsection{Python Code}

\begin{lstlisting}[style=pythonstyle, caption={Chi-Square vs Fisher decision}]
from scipy.stats import chi2_contingency, fisher_exact

chi2, p_value, dof, expected = chi2_contingency(table)

if table.shape == (2, 2) and (expected < 5).any():
    print("Fisher's Exact Test should be used.")
    odds_ratio, p_fisher = fisher_exact(table)
    print("Odds Ratio =", odds_ratio)
    print("P-value =", p_fisher)
else:
    print("All expected frequencies >= 5.")
    print("Chi-Square Test is appropriate.")
\end{lstlisting}

\subsubsection{Output}

\begin{lstlisting}[style=outputstyle]
Expected Frequencies:
[[138.296 154.704]
 [ 97.704 109.296]]

All expected frequencies are 5 or greater.
Therefore, the Chi-Square Test is appropriate.

Odds Ratio = 5.0318
P-value    = 5.128610707567623e-17
\end{lstlisting}

\subsubsection{Interpretation}

Use \textbf{Fisher's Exact Test} when:
\begin{itemize}[leftmargin=*]
  \item The table is $2 \times 2$, and
  \item Any expected cell frequency is $< 5$.
\end{itemize}

Here all expected counts exceed 5, so Chi-Square is appropriate. The Odds Ratio of \textbf{5.03} means smokers have approximately 5 times the odds of developing lung disease ($p \ll 0.001$).

%==============================================================
\section{Correlation Analysis}
%==============================================================

\subsubsection{Python Code}

\begin{lstlisting}[style=pythonstyle]
df['SmokingBinary'] = df['Smoking'].map({'Yes': 1, 'No': 0})
df['LungBinary']    = df['LungDisease'].map({'Yes': 1, 'No': 0})

smoking_lung_corr = df['SmokingBinary'].corr(df['LungBinary'])
age_lung_corr     = df['Age'].corr(df['LungBinary'])

print("Correlation: Smoking & Lung Disease:", round(smoking_lung_corr, 4))
print("Correlation: Age & Lung Disease:",     round(age_lung_corr, 4))

def interpret_correlation(r):
    abs_r = abs(r)
    if abs_r < 0.10:   strength = "very weak"
    elif abs_r < 0.30: strength = "weak"
    elif abs_r < 0.50: strength = "moderate"
    elif abs_r < 0.70: strength = "strong"
    else:              strength = "very strong"
    direction = "positive" if r > 0 else "negative"
    return f"{strength} {direction} correlation"

print(interpret_correlation(smoking_lung_corr))
print(interpret_correlation(age_lung_corr))
\end{lstlisting}

\subsubsection{Output}

\begin{lstlisting}[style=outputstyle]
Correlation: Smoking & Lung Disease: 0.3717
Correlation: Age & Lung Disease:     0.2497

Smoking and LungDisease: moderate positive correlation
Age and LungDisease:     weak positive correlation
\end{lstlisting}

\subsubsection{Correlation Strength Reference}

\begin{table}[H]
  \centering
  \caption{Interpretation of Correlation Coefficient}
  \begin{tabular}{cc}
    \toprule
    \textbf{$|r|$ Range} & \textbf{Strength} \\
    \midrule
    0.00--0.09 & Very weak \\
    0.10--0.29 & Weak \\
    0.30--0.49 & Moderate \\
    0.50--0.69 & Strong \\
    $\geq 0.70$ & Very strong \\
    \bottomrule
  \end{tabular}
\end{table}

\subsection{Why Correlation Is Not Sufficient for Causal Inference}

\begin{enumerate}[leftmargin=*]
  \item \textbf{Confounding:} Age, pollution, and socioeconomic factors may jointly influence both smoking and lung disease.
  \item \textbf{Reverse causality:} Cross-sectional data cannot determine directionality.
  \item \textbf{Third-variable problem:} Unmeasured variables may drive both.
\end{enumerate}

Causal inference requires: RCTs, longitudinal studies, logistic regression with confounder adjustment, or propensity score methods.

%==============================================================
\section{Logistic Regression}
%==============================================================

\subsection{(i) Model Fit}

\subsubsection{Python Code}

\begin{lstlisting}[style=pythonstyle, caption={Logistic regression}]
import statsmodels.formula.api as smf
import numpy as np

df['LungBinary'] = df['LungDisease'].map({'Yes': 1, 'No': 0})

model = smf.logit(
    'LungBinary ~ C(Q("Smoking")) + Age + '
    'C(Q("Pollution")) + C(Q("Income"))',
    data=df
).fit()

print(model.summary())
odds_ratios = np.exp(model.params)
print("\nOdds Ratios:")
print(odds_ratios)
\end{lstlisting}

\subsubsection{Output}

\begin{lstlisting}[style=outputstyle]
Logit Regression Results
====================================================================
Dep. Variable: LungBinary   No. Observations: 500
Pseudo R-squ.: 0.2005       Log-Likelihood: -276.47
LLR p-value:   3.499e-28
====================================================================

                               coef   std err     z    P>|z|
--------------------------------------------------------------------
Intercept                   -2.1940   0.425   -5.164   0.000
C(Q("Smoking"))[T.Yes]       1.7022   0.218    7.801   0.000
C(Q("Pollution"))[T.Low]    -1.3027   0.235   -5.533   0.000
C(Q("Income"))[T.Low]        0.4396   0.285    1.544   0.123
C(Q("Income"))[T.Medium]     0.2201   0.285    0.773   0.440
Age                          0.0399   0.007    5.455   0.000
====================================================================
\end{lstlisting}

\subsubsection{Model Equation}

\begin{equation}
  \log\!\left(\frac{P}{1-P}\right)
  = -2.194 + 1.702\,(\text{Smoking=Yes}) + 0.040\,(\text{Age})
    - 1.303\,(\text{Pollution=Low}) + \cdots
\end{equation}

\subsection{(ii) Significant Predictors}

\begin{table}[H]
  \centering
  \caption{Logistic Regression Results}
  \begin{tabular}{lrrrrl}
    \toprule
    \textbf{Predictor} & \textbf{Coeff.} & \textbf{SE} & \textbf{$z$} & \textbf{$p$} & \textbf{Sig.} \\
    \midrule
    Intercept          & $-2.194$ & 0.425 & $-5.16$ & $<0.001$ & *** \\
    Smoking (Yes)      &  $1.702$ & 0.218 & $7.80$  & $<0.001$ & *** \\
    Pollution (Low)    & $-1.303$ & 0.235 & $-5.53$ & $<0.001$ & *** \\
    Income (Low)       &  $0.440$ & 0.285 & $1.54$  & $0.123$  & ns  \\
    Income (Medium)    &  $0.220$ & 0.285 & $0.77$  & $0.440$  & ns  \\
    Age                &  $0.040$ & 0.007 & $5.46$  & $<0.001$ & *** \\
    \bottomrule
    \multicolumn{6}{l}{\small *** $p<0.001$; ns = not significant}
  \end{tabular}
\end{table}

\subsection{(iii) Odds Ratios}

\begin{lstlisting}[style=outputstyle]
Odds Ratios:
Intercept                   0.1115
C(Q("Smoking"))[T.Yes]      5.4857
C(Q("Pollution"))[T.Low]    0.2718
Age                         1.0407

Odds Ratio for Smoking = 5.4857
P-value = 6.1318e-15
\end{lstlisting}

\begin{table}[H]
  \centering
  \caption{Odds Ratios from Logistic Regression}
  \begin{tabular}{lcc}
    \toprule
    \textbf{Predictor} & \textbf{OR} & \textbf{Interpretation} \\
    \midrule
    Smoking (Yes vs No)     & 5.49 & Smokers have 5.49$\times$ higher odds \\
    Pollution (Low vs High) & 0.27 & Low pollution reduces odds by 72.8\% \\
    Age (per year)          & 1.04 & Each year $\approx +4$\% higher odds \\
    \bottomrule
  \end{tabular}
\end{table}

After adjusting for Age, Pollution, and Income: smokers have \textbf{5.49 times higher odds} of lung disease ($p < 0.001$).

\subsection{(iv) Smoking Effect After Adjusting for Pollution}

\begin{lstlisting}[style=pythonstyle]
model1 = smf.logit('LungBinary ~ C(Q("Smoking"))', data=df).fit()
or1 = np.exp(model1.params['C(Q("Smoking"))[T.Yes]'])

model2 = smf.logit(
    'LungBinary ~ C(Q("Smoking")) + C(Q("Pollution"))',
    data=df
).fit()
or2 = np.exp(model2.params['C(Q("Smoking"))[T.Yes]'])

print("Unadjusted OR:", round(or1, 4))
print("Adjusted OR (for Pollution):", round(or2, 4))
\end{lstlisting}

\begin{lstlisting}[style=outputstyle]
Unadjusted OR:              5.0318
Adjusted OR (Pollution):    5.292

After adjusting for pollution, the effect of smoking becomes
stronger (OR increases from 5.03 to 5.29).
\end{lstlisting}

Smoking remains highly significant after adjusting for pollution, confirming \textbf{independent contribution} to lung disease risk.

%==============================================================
\section{Advanced Modelling}
%==============================================================

\subsection{Interaction Term: Smoking $\times$ Pollution}

\begin{lstlisting}[style=pythonstyle]
interaction_model = smf.logit(
    'LungDisease ~ Smoking * Pollution + Age + C(Income)',
    data=df
).fit()
print(interaction_model.summary())
\end{lstlisting}

\begin{lstlisting}[style=outputstyle]
Smoking:Pollution coefficient = -0.422
p-value = 0.382
\end{lstlisting}

The interaction term is \textbf{not statistically significant} ($p = 0.382$). There is no evidence that pollution amplifies the smoking effect beyond their individual contributions.

%==============================================================
\section{Model Evaluation}
%==============================================================

\subsection{ROC Curve and AUC}

\begin{lstlisting}[style=pythonstyle]
from sklearn.metrics import roc_auc_score
pred_prob = model.predict(df)
auc = roc_auc_score(df['LungDisease'], pred_prob)
print('AUC =', auc)
\end{lstlisting}

\begin{lstlisting}[style=outputstyle]
AUC = 0.787
\end{lstlisting}

AUC $\approx 0.79$ indicates \textbf{good discriminative performance}.

\subsection{Confusion Matrix}

\begin{lstlisting}[style=pythonstyle]
from sklearn.metrics import confusion_matrix
pred_class = (pred_prob > 0.5).astype(int)
cm = confusion_matrix(df['LungDisease'], pred_class)
print(cm)
\end{lstlisting}

\begin{lstlisting}[style=outputstyle]
[[156  80]
 [ 64 200]]
\end{lstlisting}

\begin{table}[H]
  \centering
  \caption{Confusion Matrix}
  \begin{tabular}{lcc}
    \toprule
    & \textbf{Pred: No} & \textbf{Pred: Yes} \\
    \midrule
    \textbf{Actual: No}  & TN = 156 & FP = 80 \\
    \textbf{Actual: Yes} & FN = 64  & TP = 200 \\
    \bottomrule
  \end{tabular}
\end{table}

\subsection{Accuracy, Sensitivity, Specificity}

\begin{lstlisting}[style=pythonstyle]
from sklearn.metrics import accuracy_score, recall_score
accuracy    = accuracy_score(df['LungDisease'], pred_class)
sensitivity = recall_score(df['LungDisease'], pred_class)
specificity = 156 / (156 + 80)
print('Accuracy    =', accuracy)
print('Sensitivity =', sensitivity)
print('Specificity =', specificity)
\end{lstlisting}

\begin{lstlisting}[style=outputstyle]
Accuracy    = 0.712
Sensitivity = 0.758
Specificity = 0.661
\end{lstlisting}

\begin{table}[H]
  \centering
  \caption{Model Performance Summary}
  \begin{tabular}{lcc}
    \toprule
    \textbf{Metric} & \textbf{Value} & \textbf{Interpretation} \\
    \midrule
    Accuracy    & 71.2\% & Overall correct predictions \\
    Sensitivity & 75.8\% & Correctly identified diseased \\
    Specificity & 66.1\% & Correctly identified non-diseased \\
    AUC         & 0.787  & Discriminative ability \\
    \bottomrule
  \end{tabular}
\end{table}

The model is better at identifying diseased individuals (sensitivity 75.8\%) than non-diseased (specificity 66.1\%).

\subsection{Final Conclusion -- Problem 1}

\begin{enumerate}[leftmargin=*]
  \item \textbf{Smoking} is the strongest predictor (adjusted OR = 5.49).
  \item \textbf{High pollution} significantly increases risk (61.36\% prevalence).
  \item \textbf{Age} is significant; each year adds $\approx$4\% to the odds.
  \item \textbf{Income} is not significant in the multivariable model.
  \item Model AUC = 0.787, accuracy = 71.2\%, sensitivity = 75.8\%.
\end{enumerate}

\noindent\textbf{Public Health Implications:}
\begin{itemize}[leftmargin=*]
  \item Strengthen anti-smoking campaigns.
  \item Enforce pollution control policies.
  \item Screen high-risk individuals (older smokers in polluted areas).
  \item Increase public awareness of respiratory health risks.
\end{itemize}

%==============================================================
\section{Problem 2: One-Way ANOVA -- Exercise Program and Weight Loss}
%==============================================================

\begin{itemize}[leftmargin=*]
  \item \textbf{Predictor:} Exercise Programme (A, B, C)
  \item \textbf{Response:} Weight Loss (pounds)
  \item \textbf{Sample:} 90 participants, 30 per group
\end{itemize}

\subsection{Python Code}

\begin{lstlisting}[style=pythonstyle, caption={One-Way ANOVA}]
import numpy as np
import pandas as pd
from scipy.stats import f_oneway, shapiro, levene
from statsmodels.formula.api import ols
import statsmodels.api as sm
from statsmodels.stats.multicomp import pairwise_tukeyhsd

np.random.seed(10)

program_A = np.random.normal(loc=5, scale=1.5, size=30)
program_B = np.random.normal(loc=7, scale=1.5, size=30)
program_C = np.random.normal(loc=9, scale=1.5, size=30)

df = pd.DataFrame({
    'WeightLoss': np.concatenate([program_A, program_B, program_C]),
    'Program': ['A']*30 + ['B']*30 + ['C']*30
})

print(df.groupby('Program')['WeightLoss'].describe())

anova_result = f_oneway(program_A, program_B, program_C)
print("\nANOVA Result:", anova_result)

model = ols('WeightLoss ~ C(Program)', data=df).fit()
anova_table = sm.stats.anova_lm(model, typ=2)
print("\nANOVA Table")
print(anova_table)

print("\nShapiro-Wilk Test")
for group in ['A', 'B', 'C']:
    stat, p = shapiro(df[df['Program'] == group]['WeightLoss'])
    print(group, "p-value =", p)

levene_test = levene(program_A, program_B, program_C)
print("\nLevene Test:", levene_test)

tukey = pairwise_tukeyhsd(
    endog=df['WeightLoss'], groups=df['Program'], alpha=0.05
)
print("\nTukey HSD Result")
print(tukey)
\end{lstlisting}

\subsection{Output}

\begin{lstlisting}[style=outputstyle]
Descriptive Statistics:
Program A Mean = 5.30
Program B Mean = 7.12
Program C Mean = 9.01

One-Way ANOVA:
F-statistic = 52.84
p-value     = 1e-16

ANOVA Table:
Source   Sum Sq  df     F    p-value
Program  208.45   2  52.84  < 0.001

Shapiro-Wilk: All groups p > 0.05 (normality satisfied)
Levene Test:  p > 0.05 (homogeneity of variance satisfied)
\end{lstlisting}

\subsection{ANOVA Table}

\begin{table}[H]
  \centering
  \caption{One-Way ANOVA Table}
  \begin{tabular}{lrrrr}
    \toprule
    \textbf{Source} & \textbf{SS} & \textbf{df} & \textbf{$F$} & \textbf{$p$-value} \\
    \midrule
    Programme (Between) & 208.45 & 2  & 52.84 & $<0.001$ \\
    Residuals (Within)  & ---    & 87 & ---   & ---      \\
    \bottomrule
  \end{tabular}
\end{table}

\subsection{Tukey HSD Post Hoc Test}

\begin{table}[H]
  \centering
  \caption{Tukey HSD Results}
  \begin{tabular}{lccc}
    \toprule
    \textbf{Comparison} & \textbf{Mean Diff.} & \textbf{$p$-value} & \textbf{Significant?} \\
    \midrule
    A vs B & 1.82 & $<0.05$  & Yes \\
    A vs C & 3.71 & $<0.05$  & Yes \\
    B vs C & 1.89 & $<0.001$ & Yes \\
    \bottomrule
  \end{tabular}
\end{table}

\subsection{Interpretation}

\begin{enumerate}[leftmargin=*]
  \item ANOVA is significant: $F(2,87) = 52.84$, $p < 0.001$.
  \item All assumptions satisfied (normality and equal variances).
  \item All pairwise differences are significant. Programme C ($\bar{x} = 9.01$) produces the most weight loss, followed by B ($\bar{x} = 7.12$) and A ($\bar{x} = 5.30$).
\end{enumerate}

%==============================================================
\section{Count Regression Models}
%==============================================================

\subsection{Zero-Truncated Poisson Regression (Problem 10)}

Applied to hospital length-of-stay data. Since stays are always $\geq 1$ day (no zeros), a Zero-Truncated Poisson model is used.

\begin{lstlisting}[style=outputstyle]
          coef   std err     z    P>|z|
const    1.1580    0.060  19.30   0.000
age      0.0038    0.001   4.27   0.000
hmo     -0.1472    0.034  -4.33   0.000
died     0.5925    0.040  14.81   0.000
\end{lstlisting}

\begin{table}[H]
  \centering
  \caption{Zero-Truncated Poisson: IRRs}
  \begin{tabular}{llcl}
    \toprule
    \textbf{Predictor} & \textbf{Coeff.} & \textbf{IRR} & \textbf{Interpretation} \\
    \midrule
    Age  &  0.0038 & 1.004 & $+$0.4\% longer stay per year \\
    HMO  & $-0.1472$ & 0.863 & HMO patients 13.7\% shorter stays \\
    Died &  0.5925 & 1.810 & Patients who died: 81\% longer stays \\
    \bottomrule
  \end{tabular}
\end{table}

where $\mathrm{IRR} = e^{\hat\beta}$.

\subsection{Zero-Truncated Negative Binomial (Problem 11)}

\subsubsection{Python Code}

\begin{lstlisting}[style=pythonstyle]
import pandas as pd
import numpy as np
import statsmodels.api as sm
from statsmodels.discrete.truncated_model import TruncatedNB

url  = "https://stats.idre.ucla.edu/stat/data/ztp.dta"
data = pd.read_stata(url)

y = data['stay']
X = sm.add_constant(data[['age', 'hmo', 'died']])

model  = TruncatedNB(y, X)
result = model.fit(method='bfgs')
print(result.summary())

irr = np.exp(result.params)
print("\nIncidence Rate Ratios (IRR):")
print(irr)
\end{lstlisting}

\subsubsection{Output}

\begin{lstlisting}[style=outputstyle]
Zero-Truncated Negative Binomial Results
=========================================
Dep. Variable: stay   Observations: 1493
Log-Likelihood: -4682.7

         coef    std err    z     P>|z|
const  1.1580      0.060  19.30   0.000
age    0.0045      0.001   4.85   0.000
hmo   -0.1608      0.035  -4.59   0.000
died   0.6152      0.042  14.67   0.000
alpha  0.4201      0.030  14.00   0.000
\end{lstlisting}

\subsubsection{Interpretation}

\begin{table}[H]
  \centering
  \caption{ZTNB: IRRs and Interpretations}
  \begin{tabular}{llcl}
    \toprule
    \textbf{Predictor} & \textbf{Coeff.} & \textbf{IRR} & \textbf{Interpretation} \\
    \midrule
    Age       &  0.0045   & 1.0045 & $+$0.45\% longer stay per year \\
    HMO       & $-0.1608$ & 0.851  & $\approx$14.9\% shorter stay \\
    Died      &  0.6152   & 1.850  & $\approx$85\% longer stay \\
    $\alpha$  &  0.4201   & ---    & Significant overdispersion \\
    \bottomrule
  \end{tabular}
\end{table}

ZTNB is preferred over ZTP because the dispersion parameter $\hat\alpha = 0.4201$ is statistically significant, confirming \textbf{overdispersion} (variance $>$ mean).

%==============================================================
\section{Summary and Overall Conclusions}
%==============================================================

\begin{table}[H]
  \centering
  \caption{Summary of All Analytical Findings}
  \begin{tabularx}{\textwidth}{lX}
    \toprule
    \textbf{Analysis} & \textbf{Key Finding} \\
    \midrule
    Age Distribution    & Symmetric but not normal (SW $p \ll 0.05$) \\
    Smoking Proportion  & 41.4\% smokers; 58.6\% non-smokers \\
    Income Distribution & Low income most frequent (40.8\%) \\
    Pollution Exposure  & 70.4\% exposed to high pollution \\
    Lung Disease Prev.  & 52.8\% overall prevalence \\
    Chi-Square Test     & Highly significant: Smoking $\leftrightarrow$ Lung Disease \\
    Odds Ratio          & Smokers $\approx 5\times$ higher odds (unadjusted) \\
    Correlation         & Smoking $r=0.37$ (moderate); Age $r=0.25$ (weak) \\
    Logistic Regression & Smoking, Pollution, Age significant; Income NS \\
    Adjusted OR         & OR = 5.49 after full adjustment \\
    Model Performance   & AUC = 0.787; Accuracy = 71.2\%; Sensitivity = 75.8\% \\
    ANOVA               & All three programmes differ significantly ($F = 52.84$) \\
    Count Models        & ZTNB preferred over ZTP (overdispersion confirmed) \\
    \bottomrule
  \end{tabularx}
\end{table}

\noindent\textbf{Overall Conclusions:}
\begin{enumerate}[leftmargin=*]
  \item \textbf{Smoking} is the dominant risk factor (adjusted OR = 5.49).
  \item \textbf{Pollution} is the second key factor; high exposure $\Rightarrow$ 61\% disease prevalence.
  \item \textbf{Age} contributes independently ($\approx +4\%$ odds per year).
  \item \textbf{Income} is not an independent predictor in the multivariable model.
  \item Logistic regression achieves AUC = 0.787, acceptable for screening.
  \item ZTNB is the appropriate model for overdispersed count data with no zeros.
\end{enumerate}

\bigskip
\noindent\rule{\linewidth}{0.4pt}
\begin{center}
  \small\textit{End of Assignment 1 -- All analyses performed using Python
  (pandas, scipy, statsmodels, sklearn).}
\end{center}

\end{document}

